\section{Outliers and robust summaries}

An outlier is an observation that is far from the main body of the data. Outliers
matter because they can strongly affect some summaries, especially the sample
mean and standard deviation.

There is no single mathematical definition of an outlier that is correct in all
contexts. Whether an observation is unusual depends on the scale of the data, the
measurement process, and the question being asked.

\subsection*{A simple example}

Compare the datasets
\[
1,\quad 2,\quad 2,\quad 3,\quad 4
\]
and
\[
1,\quad 2,\quad 2,\quad 3,\quad 100.
\]
In the second dataset, the value \(100\) is far from the other observations. It
has a large effect on the mean:
\[
\frac{1+2+2+3+4}{5}=2.4,
\]
whereas
\[
\frac{1+2+2+3+100}{5}=21.6.
\]
The median, however, is \(2\) in both datasets.

This example shows that the mean is sensitive to extreme values, while the
median is more robust.

\subsection*{Robustness}

A summary is called robust if it is not strongly affected by a small number of
extreme observations. The median is more robust than the mean. The interquartile
range is more robust than the standard deviation.

This does not make robust summaries automatically superior. If extreme values
are genuine and important, then a sensitive summary may reveal something that a
robust summary hides. If extreme values are measurement errors, then a robust
summary may better represent the main part of the data.

\begin{center}
\begin{tabular}{lll}
\toprule
Feature & Sensitive summary & More robust summary \\
\midrule
Center & mean & median \\
Spread & standard deviation & interquartile range \\
\bottomrule
\end{tabular}
\end{center}

\subsection*{The 1.5 IQR rule}

A common descriptive convention marks observations as potential outliers if they
fall below
\[
Q_1-1.5IQR
\]
or above
\[
Q_3+1.5IQR.
\]
This rule is useful for exploratory analysis, especially in boxplots. It is not a
law of probability. It is a convention for flagging observations that deserve
attention.

\begin{remarkbox}
Calling an observation an outlier does not mean it should automatically be
removed. It means the observation should be investigated, interpreted, and
checked in context.
\end{remarkbox}

\subsection*{Outliers and modelling}

Outliers may have different explanations:
\begin{itemize}
    \item a data entry error;
    \item a measurement error;
    \item a rare but genuine event;
    \item a mixture of different populations;
    \item an inappropriate theoretical model.
\end{itemize}

For example, if most delivery times are between \(20\) and \(40\) minutes but one
recorded delivery time is \(400\) minutes, then the value might be a recording
error, a severe disruption, or evidence that the process has occasional extreme
delays. The correct interpretation cannot be decided from the number alone.

\subsection*{Relation to probability models}

Some probability models naturally produce extreme observations. Heavy-tailed
models produce large deviations more often than light-tailed models. Therefore an
observation that looks like an outlier under one model may be unsurprising under
another model.

This is another reason to distinguish data description from modelling. A boxplot
or IQR rule can flag a value as unusual in the dataset. A probability model can
then help decide whether such a value is plausible under a proposed mechanism.

\begin{summarybox}
Outliers are observations far from the main body of the data. They can strongly
affect means and standard deviations. Robust summaries such as the median and IQR
help describe the central part of the data, but unusual observations should be
interpreted in context rather than removed mechanically.
\end{summarybox}
